\chapter{总结与展望}

在完成方法构建、系统实现与正式实验之后，本文需要对整体研究工作进行归纳，并对当前工作的边界与后续发展方向作出总结。因此，本章从工作总结、不足分析与未来展望三个方面，对全文研究内容进行收束，以形成完整的论文结论。

\section{工作总结}

本文围绕工程图纸相似检索任务，完成了从数据清理、系统实现、方法设计到正式评测的完整研究流程。与旧版工作相比，本文将原先基于小规模旧测试集的实验体系重构为基于 \texttt{Dataset\_img3} 的双主线实验体系，并在 13880 张图纸、22 个类别、440 个正式查询样本上完成了新的正式实验。

在方法上，本文不再采用“清洗、局部增强、结构重排、OCR 融合”单链条递进式写法，而是构建了两条主线：
\begin{enumerate}
  \item YOLO 清洗路线；
  \item 无清洗注意力式视觉增强路线。
\end{enumerate}
在两条主线之上，本文分别考察了结构相似性增强与 OCR 文字辅助检索的作用，从而形成了 8 组正式实验模式。

实验结果表明：
\begin{enumerate}
  \item 在经过统一黑底和边框清理的大规模新数据集上，\texttt{YOLO Cleaning Only} 与 \texttt{Baseline} 指标一致，说明单独执行 YOLO 清洗已不再带来稳定收益；
  \item 结构重排序仍然是最稳定有效的增强模块；
  \item 无清洗注意力路线整体优于 YOLO 清洗路线；
  \item 当前最佳模式为 \texttt{Attention + Structure (No Cleaning)}，其 Recall@1 达到 0.8977，Recall@5 达到 0.9318，Recall@10 达到 0.9500，nDCG@10 达到 0.6882；
  \item OCR 辅助检索在两条路线中均未形成稳定正增益，说明当前标题栏字段抽取与文本融合策略仍有优化空间。
\end{enumerate}

因此，本文得到的核心结论是：对于已经完成统一预处理的工程图纸大规模数据集，\textbf{在视觉编码阶段直接增强主体结构表达，相较于继续依赖显式清洗，更适合作为工程图纸检索的主要技术路线}。

\section{不足分析}

尽管本文完成了方法重构和正式实验，但仍存在以下不足：
\begin{enumerate}
  \item 当前“注意力路线”采用的是可运行的 mask-guided pooling / attention-like 方案，尚未单独训练独立的 \texttt{spatial\_attention} 新模型；
  \item OCR 标题栏字段抽取仍然受到模板差异、分辨率和文本质量影响，因此多模态结果不够稳定；
  \item 结构描述符虽然能够改善候选重排序，但对零件图、装配图和工序图之间更高层级的语义差异建模仍不充分；
  \item 目前的多模态融合仍属于后处理式重排，尚未实现视觉与文本信息的端到端联合优化。
\end{enumerate}

\section{未来展望}

后续工作可从以下几个方向继续展开：
\begin{enumerate}
  \item 单独训练与当前路线匹配的 \texttt{spatial\_attention} 模型，进一步验证“无清洗注意力路线”在端到端训练条件下的潜力；
  \item 引入更稳定的标题栏检测与字段解析模型，提高图号、零件名称和材料字段的抽取精度；
  \item 将 OCR 融合从后处理重排升级为训练期联合约束，使文本语义信息更早参与特征学习；
  \item 引入更精细的图纸层级标签和部件级标注，增强系统对装配图、零件图和工序图之间语义层次差异的辨别能力；
  \item 继续扩充标准化工程图纸数据集，提升实验结论的泛化性与学术说服力。
\end{enumerate}

总体而言，本文围绕工程图纸相似检索问题完成了较系统的研究与实现工作，既给出了可运行的工程方案，也形成了基于大规模清理数据集的正式实验结论。虽然当前工作仍存在若干局限，但相关方法、系统与实验框架已经为后续深入研究奠定了较为扎实的基础。
